\section{Introduction to Wireless Networks (7.1)}
%=============================================================================

\subsection{Elements of a Wireless Network}

\begin{keybox}[Key Components of a Wireless Network]
A wireless network consists of several key elements:
\begin{itemize}[noitemsep]
    \item \textbf{Wireless hosts:} End-user devices (laptops, smartphones, tablets) running applications. May be stationary or mobile.
    \item \textbf{Base stations:} Fixed infrastructure devices (cell towers, WiFi access points) that relay packets between wired network and wireless hosts within coverage area.
    \item \textbf{Wireless links:} Communication medium connecting hosts to base stations (or hosts to hosts in ad hoc mode). Various link standards differ in range, data rate, and frequency band.
    \item \textbf{Infrastructure (wired) network:} The backbone connecting base stations to each other and to the global Internet.
\end{itemize}
\end{keybox}

\begin{imagebox}[Elements of a Wireless Network]
\textbf{Textbook:} Kurose \& Ross 8e, Figure 7.1\\
\textbf{Google:} ``wireless network elements infrastructure ad-hoc base station diagram''\\
\textit{Shows wireless hosts, base stations, wireless links, and the wired infrastructure backbone.}
\end{imagebox}

\subsection{Infrastructure Mode vs.\ Ad Hoc Mode}

\begin{keybox}[Infrastructure Mode vs.\ Ad Hoc Mode]
\textbf{Infrastructure mode:}
\begin{itemize}[noitemsep]
    \item Hosts connect \emph{through} a base station (access point or cell tower).
    \item Base station connects to wired network -- provides Internet access.
    \item Handoff: host moves from one base station's coverage to another (typical in cellular networks, possible in WiFi).
\end{itemize}

\medskip
\textbf{Ad hoc (infrastructure-less) mode:}
\begin{itemize}[noitemsep]
    \item No base station. Hosts communicate \emph{directly} with each other.
    \item Hosts must self-organize, self-route.
    \item Examples: Bluetooth devices, vehicular networks, emergency scenarios where infrastructure is unavailable.
\end{itemize}
\end{keybox}

\subsection{Wireless Network Taxonomy}

\begin{keybox}[Wireless Network Taxonomy -- 2x2 Matrix]
Networks are classified along two dimensions:

\begin{center}
\begin{tabular}{|l|l|l|}
\hline
 & \textbf{Single hop} & \textbf{Multiple hops} \\
\hline
\textbf{Infrastructure} & WiFi 802.11, 4G/5G LTE & Wireless mesh, relay \\
\hline
\textbf{No infrastructure} & Bluetooth, ad hoc WiFi & MANET, VANET \\
\hline
\end{tabular}
\end{center}

\begin{itemize}[noitemsep]
    \item \textbf{Single-hop, infrastructure:} Host connects to base station in one hop. Most common (home WiFi, cellular).
    \item \textbf{Single-hop, no infrastructure:} Hosts talk directly, no base station (Bluetooth).
    \item \textbf{Multi-hop, infrastructure:} Packets may relay through wireless nodes before reaching a base station (wireless mesh).
    \item \textbf{Multi-hop, no infrastructure:} Packets relay hop-by-hop with no fixed infrastructure (MANET -- Mobile Ad-hoc NETwork).
\end{itemize}
\end{keybox}

%=============================================================================
\section{Wireless Link Characteristics (7.2)}
%=============================================================================

\subsection{Challenges of the Wireless Channel}

\begin{warnbox}[Why Wireless Is Harder Than Wired]
Three fundamental challenges in wireless communications:
\begin{enumerate}[noitemsep]
    \item \textbf{Path loss / decreased signal strength:} Signal energy attenuates as it propagates through space (even in free space). Signal strength falls off roughly as $1/d^2$ (free-space) or faster with obstacles.
    \item \textbf{Interference from other sources:} ISM bands (2.4 GHz, 5 GHz used by WiFi) are unlicensed -- shared with microwaves, Bluetooth, cordless phones. Licensed cellular bands still experience inter-cell interference.
    \item \textbf{Multipath propagation:} Signal reflects off objects (buildings, ground, furniture) and arrives at receiver via multiple paths with different delays. Causes constructive/destructive interference -- the channel changes over time (fading).
\end{enumerate}
All three mean: wireless links have \emph{higher} bit error rates than wired links, and the error rate is highly variable.
\end{warnbox}

\subsection{SNR vs.\ BER -- The Fundamental Tradeoff}

\begin{formulabox}[SNR and BER -- Signal-to-Noise Ratio vs.\ Bit Error Rate]
\textbf{SNR (Signal-to-Noise Ratio):} Ratio of received signal power to noise power, expressed in dB.
\[ \text{SNR}_{dB} = 10 \cdot \log_{10}\!\left(\frac{P_\text{signal}}{P_\text{noise}}\right) \]

\textbf{BER (Bit Error Rate):} Fraction of transmitted bits received in error.

\medskip
\textbf{Key relationships:}
\begin{itemize}[noitemsep]
    \item \textbf{Higher SNR $\Rightarrow$ lower BER.} More signal power relative to noise makes it easier to decode bits correctly.
    \item \textbf{Higher bit-rate modulation $\Rightarrow$ higher BER} (for the same SNR). E.g., QAM-256 encodes more bits per symbol but requires higher SNR to achieve the same BER as BPSK.
    \item As a host moves farther from the base station, SNR decreases and BER increases.
\end{itemize}
\end{formulabox}

\begin{keybox}[Modulation and Adaptive Rate]
Different modulation schemes achieve different tradeoffs:
\begin{itemize}[noitemsep]
    \item \textbf{BPSK} (2 states): 1 bit/symbol -- robust, low BER at given SNR, but slow.
    \item \textbf{QPSK} (4 states): 2 bits/symbol.
    \item \textbf{QAM-16} (16 states): 4 bits/symbol.
    \item \textbf{QAM-256} (256 states): 8 bits/symbol -- fast, but requires high SNR.
\end{itemize}

\medskip
\textbf{Rate adaptation / adaptive modulation and coding (AMC):}\\
Transmitter and receiver dynamically adjust the modulation scheme based on measured channel quality (SNR estimate).
\begin{itemize}[noitemsep]
    \item Good channel (high SNR) $\rightarrow$ use high-rate modulation (e.g., QAM-64) to maximize throughput.
    \item Poor channel (low SNR, e.g., host far from AP) $\rightarrow$ fall back to lower-rate, more robust modulation (e.g., BPSK).
    \item Both 802.11 WiFi and 4G/5G LTE use AMC.
\end{itemize}
\end{keybox}

\begin{exambox}[Exam -- SNR/BER/Modulation - V23 Q1.4.3 and V24 Q1.4.3]
\textbf{V23 Q1.4.3 -- Which statements about wireless physical layer are true?}
\begin{itemize}[noitemsep]
    \item \textbf{a) Higher SNR means lower BER.} \checkmark\ \textbf{TRUE}
    \item b) Wireless links have lower BER than wired links. \texttimes\ FALSE -- wireless is generally worse.
    \item c) Higher modulation bit rate means lower BER at same SNR. \texttimes\ FALSE -- higher bit rate means higher BER.
    \item d) Multipath propagation eliminates interference. \texttimes\ FALSE -- it causes fading/interference.
    \item \textbf{e) As SNR decreases, using a lower modulation bit rate can maintain BER.} \checkmark\ \textbf{TRUE} -- this is exactly rate adaptation.
\end{itemize}

\medskip
\textbf{V24 Q1.4.3 -- SNR/BER/Modulation statements:}
\begin{itemize}[noitemsep]
    \item \textbf{a) Lower SNR means higher BER.} \checkmark\ \textbf{TRUE} (equivalent to: lower SNR $\Rightarrow$ worse channel $\Rightarrow$ more errors).
    \item b) Lower SNR means lower BER. \texttimes\ FALSE.
    \item c) Higher modulation rate means lower BER at given SNR. \texttimes\ FALSE.
    \item \textbf{d) Higher modulation bit rate means higher BER at given SNR.} \checkmark\ \textbf{TRUE}
    \item e) Path loss increases signal strength. \texttimes\ FALSE -- path loss \emph{decreases} signal strength.
\end{itemize}

\medskip
\textbf{Rule of thumb for exam:} Higher SNR = better; higher bit-rate modulation = more errors per bit for same SNR.
\end{exambox}

\subsection{Hidden Terminal Problem}

\begin{warnbox}[The Hidden Terminal Problem]
In wireless networks, node A may be able to hear node B, but \emph{not} hear node C which is also transmitting to B. From A's perspective, the channel appears idle even though B is experiencing a collision.

\begin{center}
\texttt{A ~~---~~> B <---~~ C} \quad (A and C cannot hear each other)
\end{center}

\begin{itemize}[noitemsep]
    \item A transmits to B. C also transmits to B simultaneously (C cannot hear A's transmission, so C thinks channel is idle).
    \item At B, both signals collide -- but A and C don't know this.
    \item This is the \textbf{hidden terminal problem}: C is ``hidden'' from A.
    \item CSMA carrier sense does \emph{not} solve this -- A hears nothing and proceeds to transmit regardless.
    \item Solution: \textbf{RTS/CTS} mechanism (see Section~\ref{sec:rtscst}).
\end{itemize}
\end{warnbox}

\begin{imagebox}[Hidden Terminal Problem]
\textbf{Textbook:} Kurose \& Ross 8e, Figure 7.3\\
\textbf{Google:} ``hidden terminal problem wireless exposed terminal diagram''\\
\textit{Illustrates two senders that cannot hear each other but cause collision at a shared receiver.}
\end{imagebox}

%=============================================================================
\section{IEEE 802.11 -- WiFi Wireless LANs (7.3)}
%=============================================================================

\subsection{802.11 Standards Overview}

\begin{keybox}[IEEE 802.11 WiFi Standards]
\begin{center}
\begin{tabular}{lllll}
\toprule
\textbf{Standard} & \textbf{Name} & \textbf{Freq.\ band} & \textbf{Max rate} & \textbf{Notes}\\
\midrule
802.11b & WiFi 1 & 2.4 GHz & 11 Mbps & First widely deployed\\
802.11a & WiFi 2 & 5 GHz & 54 Mbps & Less range than b\\
802.11g & WiFi 3 & 2.4 GHz & 54 Mbps & Backward compat with b\\
802.11n & WiFi 4 & 2.4/5 GHz & 600 Mbps & MIMO antennas\\
802.11ac & WiFi 5 & 5 GHz & 3.47 Gbps & Wide channels, MU-MIMO\\
802.11ax & WiFi 6/6E & 2.4/5/6 GHz & 14 Gbps & OFDMA, better dense use\\
802.11af & White-Fi & TV whitespace & varies & Long range\\
802.11ah & HaLow & sub-1 GHz & varies & IoT, long range, low power\\
\bottomrule
\end{tabular}
\end{center}

\medskip
\textbf{All 802.11 standards use CSMA/CA for multiple access.}\\
All operate in \textbf{BSS (Basic Service Set)} infrastructure mode or IBSS (ad hoc/independent BSS).
\end{keybox}

\subsection{BSS Architecture and Association}

\begin{keybox}[BSS Architecture -- Basic Service Set]
\begin{itemize}[noitemsep]
    \item \textbf{BSS:} A group of wireless stations and one \textbf{Access Point (AP)}. The AP connects the BSS to the wired distribution system (router/Internet).
    \item \textbf{SSID (Service Set Identifier):} The human-readable network name (e.g., ``NTNU-WiFi'').
    \item \textbf{Channel:} 802.11 operates on channels within the frequency band. In 2.4 GHz there are 11 overlapping channels (only 1, 6, 11 are non-overlapping). Adjacent APs should use non-overlapping channels.
    \item \textbf{AP beacon frames:} AP periodically broadcasts beacon frames (typically every 100 ms) containing its SSID and MAC address.
\end{itemize}
\end{keybox}

\begin{keybox}[Passive vs.\ Active Scanning]
To associate with an AP, a host uses one of two scanning methods:

\medskip
\textbf{Passive scanning:}
\begin{enumerate}[noitemsep]
    \item Host listens on each channel for beacon frames from APs.
    \item Selects AP with best signal (or user preference).
    \item Sends association request frame to chosen AP.
    \item AP sends association response frame.
\end{enumerate}

\medskip
\textbf{Active scanning:}
\begin{enumerate}[noitemsep]
    \item Host broadcasts \textbf{probe request frame} on each channel.
    \item APs that hear it reply with \textbf{probe response frame} (SSID, rates, etc.).
    \item Host selects AP and sends association request.
    \item AP sends association response.
\end{enumerate}

\textbf{After association:} Host typically runs DHCP to obtain IP address; DHCP server often co-located with AP/router.
\end{keybox}

\begin{imagebox}[802.11 WiFi Architecture]
\textbf{Textbook:} Kurose \& Ross 8e, Figure 7.6\\
\textbf{Google:} ``802.11 WiFi BSS infrastructure AP architecture diagram''\\
\textit{Shows the BSS structure with access point, associated wireless hosts, and distribution system.}
\end{imagebox}

\subsection{CSMA/CA -- Multiple Access in 802.11}

\begin{warnbox}[Why Not CSMA/CD in Wireless?]
Wired Ethernet uses \textbf{CSMA/CD} (Collision Detection): transmit, detect collision, abort, backoff, retry.

\medskip
In wireless, collision \emph{detection} is impractical:
\begin{itemize}[noitemsep]
    \item To detect a collision, the transmitter must compare its own transmitted signal with what it receives simultaneously. In wireless, the transmitted signal is much stronger than any incoming signal -- the transmitter would drown out the received signal.
    \item Hidden terminal: two senders may be transmitting to the same receiver but cannot hear each other.
\end{itemize}
Therefore, 802.11 uses \textbf{CSMA/CA} (Collision \emph{Avoidance}) -- try to \emph{avoid} collisions rather than detect them.
\end{warnbox}

\begin{keybox}[CSMA/CA Protocol -- Step by Step]
\textbf{Sender:}
\begin{enumerate}[noitemsep]
    \item Sense channel. If channel idle for \textbf{DIFS} (Distributed Inter-Frame Space) duration, transmit entire frame.
    \item If channel busy, wait until channel becomes idle. Then wait DIFS. Then choose random backoff (binary exponential backoff). Countdown the backoff timer \emph{only} while channel is idle (freeze timer when channel is busy).
    \item When backoff reaches 0, transmit entire frame.
    \item \textbf{No collision detection} -- must transmit the \emph{entire} frame.
\end{enumerate}

\medskip
\textbf{Receiver:}
\begin{enumerate}[noitemsep]
    \item If frame received correctly, wait \textbf{SIFS} (Short Inter-Frame Space -- shorter than DIFS) and send \textbf{ACK}.
\end{enumerate}

\medskip
\textbf{Sender:}
\begin{enumerate}[noitemsep,resume]
    \item If no ACK received within timeout $\rightarrow$ collision assumed $\rightarrow$ binary exponential backoff and retransmit.
\end{enumerate}

\medskip
\textbf{Key timing:} SIFS $<$ DIFS. ACK uses SIFS so it gets priority access before any new transmission (which must wait DIFS).
\end{keybox}

\begin{formulabox}[Inter-Frame Spacing Summary]
\begin{itemize}[noitemsep]
    \item \textbf{SIFS (Short IFS):} Used for ACK, CTS (high priority). Shortest wait.
    \item \textbf{DIFS (Distributed IFS):} Used for data frames and RTS. Longer than SIFS -- ensures ACK/CTS sent before new data.
    \item \textbf{EIFS (Extended IFS):} Used after a reception error (frame not decoded). Longest.
\end{itemize}
Priority: SIFS $<$ DIFS $<$ EIFS. Shorter IFS $\rightarrow$ higher access priority.
\end{formulabox}

\subsection{RTS/CTS Mechanism}
\label{sec:rtscst}

\begin{keybox}[RTS/CTS -- Request to Send / Clear to Send]
\textbf{Problem solved:} Hidden terminal problem. Even with CSMA/CA, two hidden terminals can both sense the channel as idle and transmit simultaneously, causing a collision at the AP.

\medskip
\textbf{RTS/CTS mechanism (optional in 802.11):}
\begin{enumerate}[noitemsep]
    \item Sender waits DIFS, then sends a short \textbf{RTS (Request to Send)} frame to the AP. RTS includes: duration of intended data transmission.
    \item AP broadcasts \textbf{CTS (Clear to Send)} frame (after SIFS). CTS is heard by \emph{all} stations in AP's range (including hidden terminals).
    \item CTS reserves the channel: all stations that hear CTS defer for the duration specified.
    \item Sender transmits data frame (after SIFS).
    \item AP sends ACK (after SIFS).
\end{enumerate}

\medskip
\textbf{Why it helps:} Even if A and C are hidden from each other, both hear the CTS from the AP and know to defer. The overhead of RTS/CTS is small compared to large data frames.

\medskip
\textbf{Trade-off:} RTS/CTS adds overhead (extra frames). Worth it for large frames; typically enabled with a threshold (e.g., RTS/CTS only for frames $>$ 500 bytes).

\medskip
\textbf{Important:} RTS/CTS \emph{reduces} but does \emph{not} eliminate hidden terminal collisions (two senders can still collide on the RTS itself).
\end{keybox}

\begin{exambox}[Exam -- CSMA/CA and RTS/CTS - V23 Q1.4.4 and V24 Q1.4.4]
\textbf{V23 Q1.4.4 -- Which are elements of 802.11 MAC?}
\begin{itemize}[noitemsep]
    \item \textbf{a) CSMA/CA.} \checkmark\ \textbf{TRUE} -- the core MAC protocol.
    \item b) CSMA/CD. \texttimes\ FALSE -- that is Ethernet.
    \item \textbf{c) ACK for every frame.} \checkmark\ \textbf{TRUE} -- receiver sends ACK after every successfully received frame (SIFS wait).
    \item \textbf{d) RTS/CTS mechanism.} \checkmark\ \textbf{TRUE} -- optional but part of the standard.
    \item \textbf{e) DIFS wait before transmission.} \checkmark\ \textbf{TRUE} -- mandatory part of CSMA/CA.
    \item f) Collision detection. \texttimes\ FALSE -- no collision detection in 802.11.
\end{itemize}

\medskip
\textbf{V24 Q1.4.4 -- CSMA/CA and RTS/CTS statements:}
\begin{itemize}[noitemsep]
    \item \textbf{a) Hidden terminal problem can still occur even with CSMA/CA.} \checkmark\ \textbf{TRUE} -- CSMA/CA alone does not solve hidden terminal; RTS/CTS is needed.
    \item b) In RTS/CTS, the sender broadcasts CTS to reserve the channel. \texttimes\ FALSE -- the \emph{AP} broadcasts CTS, not the sender.
    \item \textbf{b) In RTS/CTS, sender sends RTS to AP and AP responds with CTS.} \checkmark\ \textbf{TRUE} -- correct description of RTS/CTS.
    \item c) CSMA/CA detects collisions and retransmits. \texttimes\ FALSE -- it \emph{avoids} collisions; no detection.
    \item d) Backoff timer counts down while channel is busy. \texttimes\ FALSE -- timer is \emph{frozen} when channel is busy.
    \item \textbf{e) RTS/CTS improves channel access for large frames.} \checkmark\ \textbf{TRUE} -- reduces wasted bandwidth from large-frame collisions.
\end{itemize}

\medskip
\textbf{Key distinctions to remember:}
\begin{itemize}[noitemsep]
    \item \textbf{CA not CD:} Avoid, not detect.
    \item \textbf{AP broadcasts CTS}, not the sender.
    \item \textbf{Backoff timer freezes} when channel is busy (counts down only during idle).
    \item \textbf{SIFS $<$ DIFS}: ACK/CTS use SIFS for priority access.
\end{itemize}
\end{exambox}

\begin{exambox}[Exam -- S2025 Part II Q7 -- CSMA/CA Analysis]
\textbf{Scenario:} Six wireless nodes share an AP. At time $t=0$, each node has one packet queued to transmit. Analyze the transmission order and timing under 802.11 CSMA/CA.

\medskip
\textbf{Key points for solving such problems:}
\begin{enumerate}[noitemsep]
    \item At $t=0$ all nodes sense channel idle. All wait DIFS, then each draws a random backoff counter (uniformly from $[0, CW-1]$, initial contention window CW = 15 or 31 depending on standard).
    \item The node with the \emph{smallest} backoff counter transmits first.
    \item All other nodes freeze their timers when they detect the channel is busy (ongoing transmission).
    \item After transmission + ACK + DIFS, remaining nodes resume counting down from where they froze.
    \item If two nodes had the same backoff, they collide: both do not receive ACK, double the contention window, draw new random backoff.
\end{enumerate}

\medskip
\textbf{Typical exam questions:}
\begin{itemize}[noitemsep]
    \item Who transmits first given specific backoff values?
    \item When does a given node finish transmitting (add DIFS + transmission time + SIFS + ACK time)?
    \item What happens if two nodes draw the same backoff?
\end{itemize}

\medskip
\textbf{Formula for channel occupancy:}
\[ T_\text{total} = \text{DIFS} + \text{backoff} \cdot T_\text{slot} + T_\text{frame} + \text{SIFS} + T_\text{ACK} \]
\end{exambox}

\subsection{802.11 Frame Structure}

\begin{keybox}[802.11 Frame -- Four MAC Addresses]
Unlike Ethernet (2 MAC addresses), an 802.11 frame has \textbf{4 address fields}:

\begin{center}
\begin{tabular}{|l|l|l|}
\hline
\textbf{Field} & \textbf{In Infrastructure Mode} & \textbf{Direction} \\
\hline
Address 1 & MAC of receiver & Destination of this wireless hop \\
Address 2 & MAC of transmitter & Source of this wireless hop \\
Address 3 & MAC of router interface & Gateway router to/from wired network \\
Address 4 & Used only in ad hoc mode & -- \\
\hline
\end{tabular}
\end{center}

\medskip
\textbf{Example -- host H1 sending to server S1 via AP:}
\begin{itemize}[noitemsep]
    \item H1 $\rightarrow$ AP (wireless frame): Addr1 = AP MAC, Addr2 = H1 MAC, Addr3 = Router MAC
    \item AP $\rightarrow$ Router (Ethernet frame): Dest = Router MAC, Src = AP MAC
    \item (The IP datagram inside always has IP src = H1 IP, IP dst = S1 IP)
\end{itemize}

\textbf{Why Address 3?} The AP acts as a link-layer relay. It strips the 802.11 frame and creates an Ethernet frame for the wired LAN. Address 3 tells the AP \emph{where to forward} the packet on the wired side.
\end{keybox}

\begin{keybox}[802.11 Frame Control Field]
The \textbf{Frame Control} field (2 bytes) contains:
\begin{itemize}[noitemsep]
    \item \textbf{Protocol version} (2 bits): Always 0 for current 802.11.
    \item \textbf{Type} (2 bits): Management (00), Control (01), Data (10).
    \item \textbf{Subtype} (4 bits): E.g., beacon, probe request, RTS, CTS, ACK, data.
    \item \textbf{To DS / From DS} (1 bit each): Indicate direction relative to distribution system (infrastructure).
    \item \textbf{More fragments} (1 bit): Whether more fragments follow.
    \item \textbf{Retry} (1 bit): Whether this is a retransmission.
    \item \textbf{Power management} (1 bit): Sender going to sleep after this frame.
    \item \textbf{More data} (1 bit): AP has more frames buffered for this station.
    \item \textbf{WEP/Protected} (1 bit): Frame body is encrypted.
    \item \textbf{Order} (1 bit): Frames must be delivered in order.
\end{itemize}
\end{keybox}

\begin{imagebox}[802.11 Frame Format]
\textbf{Textbook:} Kurose \& Ross 8e, Figure 7.10\\
\textbf{Google:} ``802.11 WiFi frame format addresses fields diagram''\\
\textit{Shows the 802.11 frame layout including four address fields, frame control, duration, and payload.}
\end{imagebox}

\subsection{Mobility Within the Same Subnet}

\begin{keybox}[Mobility Within Same IP Subnet]
\textbf{Scenario:} Host H1 moves from BSS1 (AP1) to BSS2 (AP2). Both APs are in the \emph{same} IP subnet.

\medskip
\textbf{What changes and what doesn't:}
\begin{itemize}[noitemsep]
    \item \textbf{IP address stays the same} -- same subnet, no need for new DHCP (short timescale).
    \item \textbf{Association changes:} H1 disassociates from AP1 and associates with AP2.
    \item \textbf{Switch self-learning:} The Ethernet switch between AP1 and AP2 updates its forwarding table via self-learning when H1 sends a frame through AP2. 802.11 may also send a broadcast frame to ``wake up'' switch learning.
    \item \textbf{No routing change required.}
\end{itemize}

\textbf{Cross-subnet mobility} (different IP subnet) requires IP-level mobility support (Mobile IP, not covered in 7th ed. detail).
\end{keybox}

\begin{imagebox}[Mobility Management]
\textbf{Textbook:} Kurose \& Ross 8e, Figure 7.18\\
\textbf{Google:} ``wireless mobility handoff handover base station diagram''\\
\textit{Shows a mobile host handing off between base stations while maintaining an active connection.}
\end{imagebox}

\subsection{Rate Adaptation and Power Management}

\begin{keybox}[Rate Adaptation]
802.11 dynamically adjusts the modulation/coding scheme based on channel quality:
\begin{itemize}[noitemsep]
    \item Good channel (high SNR) $\rightarrow$ use high-rate modulation (e.g., QAM-64 5/6 coding) $\rightarrow$ higher throughput.
    \item Poor channel (low SNR, edge of coverage) $\rightarrow$ fall back to low-rate, robust modulation (e.g., BPSK 1/2) $\rightarrow$ maintain connection at lower rate.
    \item Feedback: receiver sends SNR estimates (or implicit: lack of ACK $\Rightarrow$ reduce rate).
\end{itemize}
This is the same AMC concept as in LTE (see Section~\ref{sec:4g}).
\end{keybox}

\begin{keybox}[802.11 Power Management]
Mobile devices want to save battery. 802.11 has a sleep/wake mechanism:
\begin{itemize}[noitemsep]
    \item Host sets \textbf{Power Management} bit in frame control = 1 to tell AP it will sleep.
    \item AP buffers frames destined for sleeping host.
    \item AP's \textbf{beacon frame} includes a list of APs with buffered frames (TIM -- Traffic Indication Map).
    \item Host wakes up before every beacon (every ~100 ms), checks TIM, stays awake if it has buffered frames, goes back to sleep otherwise.
    \item Result: host can sleep for $\sim$99\% of time if no traffic (significant power savings).
\end{itemize}
\end{keybox}

%=============================================================================
\section{Cellular Networks -- 4G and 5G (7.4)}
%=============================================================================
\label{sec:4g}

\subsection{4G LTE Architecture}

\begin{keybox}[4G LTE Network Architecture]
4G LTE (Long-Term Evolution) architecture separates the \textbf{user plane} (data) and \textbf{control plane} (signaling):

\medskip
\textbf{Radio Access Network (RAN):}
\begin{itemize}[noitemsep]
    \item \textbf{UE (User Equipment):} The mobile device (phone, tablet). Contains a SIM card with a 64-bit \textbf{IMSI (International Mobile Subscriber Identity)} that uniquely identifies the subscriber.
    \item \textbf{eNode-B (eNB):} The 4G base station. Handles radio access (scheduling, modulation, etc.) and some control functions. In 5G, called \textbf{gNode-B (gNB)}.
\end{itemize}

\medskip
\textbf{Evolved Packet Core (EPC):}
\begin{itemize}[noitemsep]
    \item \textbf{MME (Mobility Management Entity):} Control plane. Handles authentication, UE tracking, paging, handover decisions. Contacts HSS to verify IMSI.
    \item \textbf{HSS (Home Subscriber Service):} Database of subscriber information (IMSI, service profile, authentication keys).
    \item \textbf{S-GW (Serving Gateway):} User-plane gateway in the visited network. Forwards data between eNB and P-GW. Handles mobility across eNBs.
    \item \textbf{P-GW (PDN Gateway):} User-plane gateway to the Internet (Packet Data Network). Performs NAT, allocates UE IP address, enforces QoS.
\end{itemize}
\end{keybox}

\begin{imagebox}[4G/LTE Network Architecture]
\textbf{Textbook:} Kurose \& Ross 8e, Figure 7.16\\
\textbf{Google:} ``4G LTE network architecture eNodeB MME diagram''\\
\textit{Shows UE, eNode-B, S-GW, P-GW, MME, and HSS with control and user plane separation.}
\end{imagebox}

\begin{tikzpicture}[
    node distance=1.4cm and 2cm,
    box/.style={draw, rounded corners, minimum width=2.2cm, minimum height=0.8cm, align=center, font=\small},
    arr/.style={-Stealth, thick}
]
\node[box, fill=lightblue] (UE) {UE\\(Phone)};
\node[box, fill=lightblue, right=of UE] (eNB) {eNode-B\\(Base Station)};
\node[box, fill=lightyellow, right=of eNB] (SGW) {S-GW};
\node[box, fill=lightyellow, right=of SGW] (PGW) {P-GW};
\node[box, fill=lightred, above=of SGW] (MME) {MME};
\node[box, fill=lightred, above=of PGW] (HSS) {HSS};
\node[box, fill=lightgreen, right=of PGW] (Internet) {Internet};

\draw[arr] (UE) -- node[above,font=\tiny]{radio} (eNB);
\draw[arr, <->] (eNB) -- node[above,font=\tiny]{S1-U} (SGW);
\draw[arr, <->] (SGW) -- node[above,font=\tiny]{S5} (PGW);
\draw[arr, <->] (PGW) -- (Internet);
\draw[arr, <->] (MME) -- node[right,font=\tiny]{S11} (SGW);
\draw[arr, <->] (MME) -- (HSS);
\draw[arr, <->] (MME) -- node[left,font=\tiny]{S1-MME} (eNB);

\node[font=\tiny, below=0.1cm of UE] {User plane};
\node[font=\tiny, below=0.1cm of SGW] {User plane};
\end{tikzpicture}

\subsection{LTE Data Plane -- Tunneling}

\begin{keybox}[LTE Tunneling -- GTP-U Protocol]
In LTE, user data is \textbf{tunneled} from eNode-B all the way to P-GW using \textbf{GTP-U (GPRS Tunneling Protocol -- User plane)}:

\medskip
\textbf{Tunnel path:} UE $\leftrightarrow$ eNB $\xrightarrow{\text{GTP tunnel}}$ S-GW $\xrightarrow{\text{GTP tunnel}}$ P-GW $\leftrightarrow$ Internet

\begin{itemize}[noitemsep]
    \item Each user datagram is encapsulated in a GTP-U header (which adds a Tunnel Endpoint ID -- TEID).
    \item eNB $\rightarrow$ S-GW: one GTP tunnel per UE.
    \item S-GW $\rightarrow$ P-GW: another GTP tunnel per UE.
    \item The tunnels allow user traffic to flow transparently regardless of which eNB the UE is attached to.
\end{itemize}

\medskip
\textbf{Why tunneling?} Hides radio-network mobility from the IP core. When a UE moves between eNBs (handover), only the tunnel endpoints change -- the P-GW IP and the UE's allocated IP remain constant.
\end{keybox}

\subsection{LTE Link Layer}

\begin{keybox}[LTE Radio Access -- OFDM and Link-Layer Protocols]
\textbf{Physical layer:} LTE uses \textbf{OFDM (Orthogonal Frequency Division Multiplexing)}:
\begin{itemize}[noitemsep]
    \item Frequency band divided into many narrow orthogonal subcarriers.
    \item Combined FDM (across subcarriers) + TDM (time slots) = \textbf{OFDMA} in downlink.
    \item Each 1 ms subframe is divided into resource blocks (time $\times$ frequency) allocated to UEs by the base station scheduler.
    \item Each UE's resource block can use a different modulation/coding scheme (adaptive modulation per UE).
\end{itemize}

\medskip
\textbf{LTE link layer sub-layers} (between UE and eNB):
\begin{enumerate}[noitemsep]
    \item \textbf{PDCP (Packet Data Convergence Protocol):} IP header compression, ciphering (encryption).
    \item \textbf{RLC (Radio Link Control):} Fragmentation and reassembly of large packets. Reliable data transfer (ARQ -- Automatic Repeat reQuest).
    \item \textbf{MAC (Medium Access Control):} Scheduling, HARQ (Hybrid ARQ -- combines FEC + ARQ).
\end{enumerate}
\end{keybox}

\subsection{BS Association and Sleep Modes}

\begin{keybox}[BS Association in LTE]
How a UE connects to an eNB:
\begin{enumerate}[noitemsep]
    \item eNBs broadcast \textbf{synchronization signals} every 5 ms on a set of well-known frequencies.
    \item UE scans frequencies, detects synch signals, measures signal strength from multiple eNBs.
    \item UE selects the best eNB (strongest signal / highest SINR).
    \item UE sends \textbf{Random Access} message to eNB (contention-based RACH or contention-free).
    \item eNB responds with timing alignment and temporary ID.
    \item Authentication via MME and HSS: IMSI sent, HSS verifies subscriber, session established.
    \item P-GW assigns UE an IP address (via MME $\rightarrow$ S-GW $\rightarrow$ P-GW signaling).
\end{enumerate}
\end{keybox}

\begin{keybox}[LTE Sleep Modes -- DRX]
To save battery, LTE uses \textbf{DRX (Discontinuous Reception)} with two sleep modes:
\begin{itemize}[noitemsep]
    \item \textbf{Light sleep:} Wake up every 100s of milliseconds to check for downlink activity. Used when device is active but idle momentarily.
    \item \textbf{Deep sleep:} Sleep for 5--10 seconds. Used for longer idle periods (e.g., app in background). Device must re-establish connection more fully on wakeup.
\end{itemize}
\textbf{Tradeoff:} Deeper sleep $\rightarrow$ more battery savings, but longer latency to resume data transfer.
\end{keybox}

\subsection{4G vs.\ 5G}

\begin{keybox}[4G vs.\ 5G -- Key Differences]
\begin{center}
\begin{tabular}{lll}
\toprule
\textbf{Feature} & \textbf{4G LTE} & \textbf{5G NR} \\
\midrule
Peak downlink rate & $\sim$1 Gbps & Up to 20 Gbps \\
Latency & $\sim$30--50 ms & $<$1 ms (target) \\
Frequency bands & Sub-6 GHz & Sub-6 GHz + mmWave \\
Mobility & Full support & Full support \\
Architecture & Partial C/U-plane separation & Complete C/U-plane separation \\
Base station & eNode-B & gNode-B (gNB) \\
Core network & EPC & 5GC (5G Core) \\
\bottomrule
\end{tabular}
\end{center}

\medskip
\textbf{Key facts for exams:}
\begin{itemize}[noitemsep]
    \item \textbf{Both 4G and 5G support mobility / handover.}
    \item \textbf{5G has higher peak bitrate} than 4G.
    \item 5G is designed for \emph{complete} control/user-plane separation (4G has partial separation).
    \item 5G mmWave (millimeter wave, 24--100 GHz) offers very high rates but short range and limited penetration.
    \item 5G sub-6 GHz offers better range and coverage.
\end{itemize}
\end{keybox}

\begin{exambox}[Exam -- 4G/5G Statements - V23 Q1.4.5]
\textbf{V23 Q1.4.5 -- Which statements about 4G and 5G are true?}
\begin{itemize}[noitemsep]
    \item \textbf{a) 5G has higher peak bitrate than 4G.} \checkmark\ \textbf{TRUE} -- up to 20 Gbps vs.\ $\sim$1 Gbps.
    \item b) Only 5G supports mobility/handover; 4G does not. \texttimes\ FALSE -- both support mobility.
    \item \textbf{c) Both 4G and 5G support host mobility.} \checkmark\ \textbf{TRUE}
    \item d) 4G uses complete control/user-plane separation but 5G does not. \texttimes\ FALSE -- 5G is the one designed for \emph{complete} separation.
    \item e) 5G uses CSMA/CA for medium access. \texttimes\ FALSE -- that is 802.11 WiFi. LTE/5G use OFDMA and scheduled access.
\end{itemize}

\medskip
\textbf{Remember:} 5G = higher rate + lower latency + complete C/U-plane separation + both support mobility.
\end{exambox}

%=============================================================================
\section{Putting It All Together -- Wireless Summary}
%=============================================================================

\begin{keybox}[Chapter 7 -- Big Picture Summary]
\textbf{Why is wireless hard?} Three challenges: path loss, interference, multipath. These cause higher and variable BER compared to wired links.

\medskip
\textbf{SNR/BER tradeoff:} Higher SNR $\Rightarrow$ lower BER. Higher-order modulation $\Rightarrow$ higher BER at same SNR but higher throughput. Rate adaptation dynamically picks the best modulation for current channel conditions.

\medskip
\textbf{802.11 MAC:} CSMA/CA (not CD). DIFS before transmission. Random backoff (timer freezes during busy channel). Full-frame transmission. ACK after SIFS. RTS/CTS for hidden terminal mitigation.

\medskip
\textbf{802.11 frames:} 4 MAC addresses (Addr1=receiver, Addr2=transmitter, Addr3=router, Addr4=ad hoc). Frame control field encodes type, subtype, power management, WEP, To/From DS bits.

\medskip
\textbf{4G LTE:} UE $\leftrightarrow$ eNB (radio) $\leftrightarrow$ S-GW $\leftrightarrow$ P-GW $\leftrightarrow$ Internet (data plane). MME handles control (auth, mobility). HSS stores subscriber data. IMSI identifies subscriber on SIM. GTP-U tunneling end-to-end. OFDM/OFDMA radio access.

\medskip
\textbf{5G vs.\ 4G:} Higher peak rate, lower latency, complete C/U-plane separation. Both support mobility.
\end{keybox}

\begin{warnbox}[Common Exam Traps -- Chapter 7]
\begin{itemize}[noitemsep]
    \item \textbf{CSMA/CD vs.\ CSMA/CA:} 802.11 uses CA (avoidance), \emph{not} CD (detection). Wired Ethernet uses CD.
    \item \textbf{Who sends CTS?} The \emph{AP} sends CTS in response to sender's RTS -- not the sender.
    \item \textbf{Backoff timer behavior:} Timer counts down only when channel is \emph{idle}. It \emph{freezes} (not resets) when channel goes busy.
    \item \textbf{Higher SNR = better:} Do not confuse direction. Lower SNR $\Rightarrow$ higher BER $\Rightarrow$ worse performance.
    \item \textbf{Higher modulation = more fragile:} QAM-256 is faster but needs good channel. BPSK is slow but robust.
    \item \textbf{Both 4G and 5G have mobility:} A common distractor says only 5G supports mobility -- false.
    \item \textbf{5G does not use CSMA/CA:} CSMA/CA is 802.11 WiFi. LTE and 5G use OFDMA with base-station scheduling.
    \item \textbf{RTS/CTS reduces but doesn't eliminate hidden terminal:} Two nodes can still collide on their RTS frames.
\end{itemize}
\end{warnbox}

\begin{formulabox}[Quick Reference -- Key Numbers]
\begin{itemize}[noitemsep]
    \item 802.11b: 2.4 GHz, 11 Mbps
    \item 802.11g: 2.4 GHz, 54 Mbps
    \item 802.11n (WiFi 4): 2.4/5 GHz, 600 Mbps (MIMO)
    \item 802.11ac (WiFi 5): 5 GHz, 3.47 Gbps (MU-MIMO)
    \item 802.11ax (WiFi 6): 2.4/5/6 GHz, 14 Gbps (OFDMA)
    \item 802.11 beacon interval: $\sim$100 ms
    \item LTE synch signal: every 5 ms
    \item LTE light sleep: 100s of ms; deep sleep: 5--10 s
    \item LTE IMSI: 64 bits on SIM card
    \item 4G peak: $\sim$1 Gbps; 5G peak: $\sim$20 Gbps
\end{itemize}
\end{formulabox}
